% Original analytic teaching diagram, 2026-11-12.
% Fixed-length block coding by typical sets: all 4^n blocks (outer box), the typical set with about
% 2^{nH} elements and probability close to 1 (blue), each typical block gets an index of about nH bits;
% blocks outside the typical set are all sent to one reserved index and decoded wrongly (red).
\definecolor{VB}{HTML}{7188A3}\definecolor{VBF}{HTML}{E5EAF0}
\definecolor{VG}{HTML}{829C8B}\definecolor{VGF}{HTML}{E8EEE9}
\definecolor{VY}{HTML}{B4A16F}\definecolor{VYF}{HTML}{F1ECDD}
\definecolor{VR}{HTML}{A7787D}\definecolor{VRF}{HTML}{F0E5E6}
\definecolor{VP}{HTML}{9588A5}\definecolor{VPF}{HTML}{ECE8F0}
\definecolor{VI}{HTML}{3B4252}\definecolor{VM}{HTML}{6B7280}\definecolor{VL}{HTML}{4C4D4F}
\begin{tikzpicture}[x=1mm,y=1mm,>=Stealth,font=\figurefont\fontsize{8}{10}\selectfont,
lab/.style={text=VM,font=\figurefont\fontsize{7.5}{9}\selectfont},
arr/.style={->,VL,line width=.6pt}]
% all blocks
\draw[VL,line width=.7pt,rounded corners=1.5mm] (0,0) rectangle (52,36);
\node[text=VI,anchor=north west] at (1.5,35) {$\mathcal{X}^{n}$：$4^{n}$个块};
% typical set
\fill[VBF] (26,15) ellipse (16 and 9);
\draw[VB,line width=.8pt] (26,15) ellipse (16 and 9);
\node[text=VI] at (26,17.5) {$A_{\varepsilon}^{(n)}$};
\node[lab] at (26,12) {约$2^{nH}$个};
% atypical samples
\foreach \x/\y in {5/5,8/24,4/14,28/30,49/14,47/6,20/3,36/3}{
  \fill[VR] (\x,\y) circle (.7);
}
% index table
\begin{scope}[shift={(66,0)}]
\foreach \i/\t in {0/{00\cdots00},1/{00\cdots01},2/{00\cdots10},4/{11\cdots11}}{
  \draw[VG,line width=.6pt,fill=VGF] (0,30-\i*6) rectangle (26,34-\i*6);
  \node[text=VI,font=\figurefont\fontsize{7}{8}\selectfont] at (13,32-\i*6) {$\t$};
}
\node[text=VI] at (13,14.2) {$\vdots$};
\draw[VR,line width=.7pt,fill=VRF] (0,2) rectangle (26,6);
\node[text=VI,font=\figurefont\fontsize{7}{8}\selectfont] at (13,4) {保留索引};
\node[lab] at (13,39) {$\lceil n(H+\varepsilon)\rceil+1$位};
\end{scope}
\draw[arr] (36,20) .. controls (50,30) and (56,32) .. (65.5,32);
\draw[arr] (40,13) .. controls (52,14) and (56,20) .. (65.5,20);
\draw[arr,VR] (47.6,6) -- (65.5,4);
\end{tikzpicture}
